\uuid{J86c}
\titre{Transformée de Fourier d'un signal modulé et translaté}

\niveau{M1}
\module{Analyse}
\chapitre{Transformation de Fourier}
\sousChapitre{Produit, convolution, modulation et translation}
\theme{Transformée de $\cos(3t)e^{-|t-t_0|}$}
\auteur{Q. Liard}
\datecreate{ 2026-09-03}
\organisation{CentraleSupélec}
\difficulte{4}

\contenu{

\texte{
Soit $t_0\in\mathbb{R}$. On considère la fonction
$$
f(t)=\cos(3t)e^{-|t-t_0|}.
$$
On utilise la convention
$$
\widehat{u}(\xi)=\int_{\mathbb{R}}u(t)e^{-i\xi t}\,dt.
$$
Lorsque la formule d'inversion est applicable, elle s'écrit
$$
u(t)=\frac{1}{2\pi}\int_{\mathbb{R}}\widehat{u}(\xi)e^{i\xi t}\,d\xi.
$$
}

\begin{enumerate}

\item \question{Montrer que $f\in L^1(\mathbb{R})$.}

\reponse{
Pour tout $t\in\mathbb{R}$,
$$
|f(t)|=|\cos(3t)|e^{-|t-t_0|}\le e^{-|t-t_0|}.
$$
Or
$$
\int_{\mathbb{R}}e^{-|t-t_0|}\,dt
=
\int_{\mathbb{R}}e^{-|s|}\,ds
=2.
$$
Par comparaison,
$$
\boxed{f\in L^1(\mathbb{R})}.
$$
}

\item \question{Montrer que
$$
\widehat f
=
\frac{1}{2\pi}\,\widehat{\cos(3\,\cdot)}\ast\widehat{e^{-|\,\cdot-t_0|}}.
$$}

\reponse{
La fonction $t\mapsto\cos(3t)$ définit une distribution tempérée et $t\mapsto e^{-|t-t_0|}$ appartient à $L^1(\mathbb{R})\cap\mathcal{S}'(\mathbb{R})$. Comme
$$
f=\cos(3\,\cdot)\,e^{-|\,\cdot-t_0|},
$$
avec cette convention, la formule de Fourier transformant un produit en convolution donne
$$
\boxed{
\widehat f
=
\frac{1}{2\pi}\,\widehat{\cos(3\,\cdot)}\ast\widehat{e^{-|\,\cdot-t_0|}}
}.
$$
}

\item \question{Montrer que
$$
\widehat{\cos(3\,\cdot)}
=
\pi\left(\delta_{3}+\delta_{-3}\right).
$$}

\reponse{
On écrit
$$
\cos(3t)=\frac12\left(e^{3it}+e^{-3it}\right).
$$
Avec la convention de Fourier choisie,
$$
\widehat{e^{i\alpha t}}=2\pi\,\delta_{\alpha}.
$$
Par linéarité,
$$
\widehat{\cos(3\,\cdot)}
=
\pi\left(\delta_{3}+\delta_{-3}\right).
$$
}

\item \question{En déduire l'expression de la transformée de Fourier de $f$.}

\reponse{
Posons $g(t)=e^{-|t-t_0|}$. Comme $g$ est la translation de $t\mapsto e^{-|t|}$ par $t_0$,
$$
\widehat g(\xi)=e^{-i t_0\xi}\widehat{e^{-|\,\cdot|}}(\xi).
$$
Or
$$
\widehat{e^{-|\,\cdot|}}(\xi)
=
\int_{-\infty}^{+\infty}e^{-|t|}e^{-i\xi t}\,dt
=
\frac{2}{1+\xi^2}.
$$
Ainsi,
$$
\widehat g(\xi)=\frac{2e^{-i t_0\xi}}{1+\xi^2}.
$$
D'après les questions précédentes,
$$
\widehat f(\xi)
=
\frac12\left(
\widehat g\left(\xi-3\right)
+
\widehat g\left(\xi+3\right)
\right).
$$
On obtient finalement
$$
\boxed{
\widehat f(\xi)
=
\frac{e^{-i t_0(\xi-3)}}
{1+\left(\xi-3\right)^2}
+
\frac{e^{-i t_0(\xi+3)}}
{1+\left(\xi+3\right)^2}
}.
$$
}

\end{enumerate}

}
